\chapter{勒贝格积分}\label{ch:b3:lebesgue}

黎曼积分将\emph{定义域}切成许多小
区间；勒贝格积分切分的是\emph{值域}：为了积分
$f$，只需求集合 $\{f > t\}$ 的测度。这一改动看似平淡无奇，
实则具有革命性。在黎曼理论中，极限与积分始终“争执不休”，
要求一致收敛（条件很苛刻！），
借三个收敛定理得以调和 --- 单调收敛定理、
法图引理与控制收敛定理 --- 其条件
弱得几乎令人尴尬。本章构造
任意测度空间 $(X, \mathcal A,
\mu)$ 上的积分，证明三个定理，并
精确刻画它与黎曼积分的关系（有界函数
黎曼可积当且仅当它几乎处处连续），并
将含参变量积分的求导系统化 ---
这一技巧正是周末问题用来
计算 $\int_0^\infty\frac{\sin x}x\,\dd x$ 和
$\int_\R \eu^{-x^2}\dd x$ 的值。

\section{可测函数}

\begin{definition}\label{def:b3:lebesgue:measurable}
设 $(X, \mathcal A)$、$(Y, \mathcal B)$ 为可测空间。
称 $f \colon X \to Y$ 为\emph{可测函数}\index{可测
函数}，若 $f^{-1}(B) \in \mathcal A$ 对每个 $B \in
\mathcal B$ 都成立。对实值（或 $[-\infty,+\infty]$ 值的）
函数，此时 $Y = \R$ 配以其博雷尔 $\sigma$-代数，且只需
验证 $f^{-1}(\intoo t{+\infty}) = \{f > t\} \in
\mathcal A$ 对一切 $t \in \R$ 成立：满足条件的集合 $\{B :
f^{-1}(B) \in \mathcal A\}$ 构成一个 $\sigma$-代数
（原像与集合运算可交换），并包含
生成的射线（\cref{def:b3:measure:borel}，
\cref{met:b3:measure:goodsets}）。
\end{definition}

\begin{proposition}\label{prop:b3:lebesgue:stability}
(a) 可测映射的复合仍可测；连续
映射都是博雷尔可测的。
(b) 若 $f, g \colon X \to \R$ 可测，则下列函数亦可测：$f + g$，
$fg$, $\max(f,g)$, $\abs f$, $\lambda f$.
(c) 若可测函数列 $(f_n)$ 取值于
$[-\infty, +\infty]$ 中，则 $\sup_nf_n$、$\inf_nf_n$、
$\limsup f_n$、$\liminf f_n$ 皆可测；若 $f_n \to f$
逐点，则 $f$ 可测。
\end{proposition}

\begin{proof}
(a) $(g\circ f)^{-1}(B) = f^{-1}(g^{-1}(B))$；连续性
借助生成开集给出可测性
（\cref{pb:b3:measure:1}，第 10 题的一般形式）。
(b) 映射 $(f, g) \colon X \to \R^2$ 关于博雷尔
$\sigma$-代数（$\R^2$ 上的）可测 --- 只需在开方块上验证，这些方块
生成（$\R^2$ 的开集是有理
方块的可数并）：$(f,g)^{-1}(U\times V) = f^{-1}(U)\cap g^{-1}(V)$ ---
而 $+, \times, \max$ 是 $\R^2 \to \R$ 上的连续映射：复合即得。
(c) $\{\sup f_n > t\} = \bigcup_n\{f_n > t\}$; $\inf = -\sup(-)$;
$\limsup = \inf_N\sup_{n \geq N}$；逐点极限即其自身的
$\limsup$。
\end{proof}

\begin{definition}\label{def:b3:lebesgue:simple}
称取有限个值的可测函数为\emph{简单函数}\index{简单函数}，
即 $s = \sum_{i=1}^n
c_i\,\mathbf 1_{A_i}$，其中 $A_i \in \mathcal A$ 两两不交，$c_i \geq
0$（对非负情形）。其积分为
\[
\int s\,\dd\mu = \sum_i c_i\,\mu(A_i) \in [0, +\infty]
\]
（约定 $0\cdot\infty = 0$）；该值不依赖于
表示的选取（将两个划分取公共加细即可）。
\end{definition}

\begin{theorem}[简单函数的逼近]
\label{thm:b3:lebesgue:approximation}
每个可测函数 $f \colon X \to [0, +\infty]$ 都是
一列\emph{递增}的简单函数的逐点极限：
\[
s_n = \sum_{k=1}^{n2^n} \frac{k-1}{2^n}\,
\mathbf 1_{\{\frac{k-1}{2^n} \leq f < \frac k{2^n}\}}
+ n\,\mathbf 1_{\{f \geq n\}} \nearrow f .
\]
\end{theorem}

\begin{proof}
每个 $s_n$ 都是简单函数（其中各集合均为博雷尔集的原像）。
单调性：由 $n$ 过渡到 $n+1$ 时，每一二进
层各一分为二，且从不减小所赋的值（一点
满足 $\frac{k-1}{2^n} \leq f(x) < \frac k{2^n}$ 的点，取值为
$\frac{2k-2}{2^{n+1}}$ 或 $\frac{2k-1}{2^{n+1}}$，二者均 $\geq
\frac{k-1}{2^n}$；上界 $n$ 亦随之增大）。收敛性：若
$f(x) < \infty$，则对 $n > f(x)$ 有 $f(x) - s_n(x) \leq
2^{-n}$；若 $f(x) = \infty$，则 $s_n(x) = n \to \infty$。
\end{proof}

\section{积分与收敛定理}

\begin{definition}\label{def:b3:lebesgue:integral}
对可测函数 $f \geq 0$，令
\[
\int f \,\dd\mu = \sup\Bigl\{\int s\,\dd\mu : s \text{ 为简单函数},
\ 0 \leq s \leq f\Bigr\} \in [0, +\infty].
\]
由构造知其关于 $f$ 单调，且推广了简单
函数的情形（对简单 $f$，上确界在 $f$ 处取得：简单积分之比较
借助公共加细完成）。
\end{definition}

\begin{theorem}[单调收敛定理，贝波·莱维]
\label{thm:b3:lebesgue:mct}
若 $0 \leq f_n \nearrow f$ 逐点成立（可测），则
\[
\int f_n\,\dd\mu \nearrow \int f\,\dd\mu .
\]
\index{单调收敛定理}
\end{theorem}

\begin{proof}
$f$ 可测（\cref{prop:b3:lebesgue:stability}(c)）且
$\int f_n$ 增至某个 $L \leq \int f$（由单调性）。
反之，取简单函数 $s = \sum c_i\mathbf 1_{A_i} \leq f$
及 $\theta \in (0,1)$；集合 $E_n = \{f_n \geq \theta s\}$
均可测且增至 $X$（当 $s(x) > 0$ 时：$f(x)
\geq s(x) > \theta s(x)$，故最终有 $f_n(x) \geq \theta
s(x)$；当 $s(x) = 0$ 时显然成立）。于是
\[
\int f_n \geq \int_{E_n}\theta s\,\dd\mu
= \theta\sum_i c_i\,\mu(A_i \cap E_n)
\xrightarrow[n\to\infty]{} \theta\sum_ic_i\,\mu(A_i)
= \theta\int s
\]
由测度的下连续性（\cref{prop:b3:measure:basics}(c)）。
故 $L \geq \theta\int s$ 对一切 $\theta < 1$ 及一切简单函数 $s
\leq f$ 成立，于是 $L \geq \int f$。
\end{proof}

\begin{corollary}\label{cor:b3:lebesgue:additivity}
对可测的 $f, g \geq 0$ 及 $c \geq 0$：$\int(f + g) =
\int f + \int g$ 及 $\int cf = c\int f$；对由
非负可测函数构成的级数，$\int\sum_nf_n =
\sum_n\int f_n$。
\end{corollary}

\begin{proof}
对简单函数，可加性是在公共
加细上的计算。一般地，取 $s_n \nearrow f$，$t_n \nearrow g$
（见\cref{thm:b3:lebesgue:approximation}）：$s_n + t_n \nearrow f +
g$，并由单调收敛定理把可加性传到极限。级数的
结论即对部分和应用单调收敛定理。
\end{proof}

\begin{theorem}[法图引理]\label{thm:b3:lebesgue:fatou}
对可测函数 $f_n \geq 0$：
\[
\int \liminf_n f_n \,\dd\mu \;\leq\; \liminf_n \int
f_n\,\dd\mu .
\]
\index{法图引理}
\end{theorem}

\begin{proof}
记 $g_N = \inf_{n\geq N}f_n$：可测，$0 \leq g_N \nearrow
\liminf f_n$，且 $g_N \leq f_n$ 对每个 $n \geq N$ 成立，于是
$\int g_N \leq \inf_{n \geq N}\int f_n$。对左端
用单调收敛定理：$\int\liminf f_n = \lim_N\int g_N \leq
\lim_N\inf_{n\geq N}\int f_n = \liminf\int f_n$.
\end{proof}

\begin{definition}\label{def:b3:lebesgue:l1}
称可测函数 $f \colon X \to \R$（或 $\C$）为
\emph{可积}\index{可积}，若 $\int\abs
f\,\dd\mu < \infty$；此时 $\int f = \int f^+ - \int f^-$
（正部与负部；复数情形下取实部与虚部
）。积分在可积函数上是线性的
（分解并重组正部；复数情形化归为实数情形
），且满足 $\abs{\int f} \leq
\int\abs f$（实数情形：$\pm\int f = \int(\pm f) \leq
\int\abs f$；复数情形：乘以模为 1 的常数
使积分为实数）。若某性质\emph{几乎
处处}（a.e.）\index{几乎处处}成立，即它只在
一个 $\mu$ 零测集上失效；在零测集上修改 $f$ 不改变
任何积分（两者之差受 $\infty\cdot\mathbf
1_N$ 控制，其积分为 $0$）。
\end{definition}

\begin{theorem}[控制收敛定理]\label{thm:b3:lebesgue:dct}
设 $f_n \to f$ 几乎处处，且 $\abs{f_n} \leq g$ 几乎处处，对固定的
\emph{可积}函数 $g$。则 $f$ 可积，且
\[
\int f_n\,\dd\mu \longrightarrow \int f\,\dd\mu,
\qquad\text{即}\quad \int\abs{f_n - f}\,\dd\mu \to 0 .
\]
\index{控制收敛定理}
\end{theorem}

\begin{proof}
去掉一个零测集，使假设逐点成立。$\abs f
\leq g$，故 $f$ 可积。函数 $h_n = 2g - \abs{f_n
- f} \geq 0$ 满足 $\liminf h_n = 2g$；由法图引理得
\[
\int 2g \leq \liminf\int\bigl(2g - \abs{f_n - f}\bigr)
= \int 2g - \limsup\int\abs{f_n - f},
\]
于是 $\limsup\int\abs{f_n - f} \leq 0$（减法是合法的：
$\int 2g < \infty$）。最后 $\abs{\int f_n - \int f} \leq
\int\abs{f_n - f} \to 0$.
\end{proof}

\begin{method}\label{met:b3:lebesgue:convergence}
面对 $\lim_n\int f_n$，依次尝试 --- (1) 该序列是否单调
（或非负项级数）？用单调收敛定理，
无需可积性。(2) 是否存在一个可积的
控制函数 $g \geq \abs{f_n}$，可由粗略估计求得
（“$\sup_n$”式估计）？用控制收敛定理。(3) 若既无控制，
又不单调？法图引理仍给出一侧界，而等号可能确实不成立，
例如逃逸的尖峰 $f_n = n\mathbf
1_{\intoo0{1/n}}$ 满足 $\int f_n = 1$，但 $f_n \to 0$ 几乎处处。
控制恰好阻止质量逃逸到
无穷，无论沿竖直方向还是水平方向。
\end{method}

\section{黎曼与勒贝格}

\begin{theorem}[勒贝格准则]
\label{thm:b3:lebesgue:riemann}
设 $f \colon \intcc ab \to \R$ 为\emph{有界}函数。则 $f$
黎曼可积，当且仅当 $f$ 关于 $\lambda$ 几乎
处处连续；此时 $f$ 勒贝格可积，且两个
积分相等。\index{黎曼积分与勒贝格积分}
\end{theorem}

\begin{proof}
对区间的一个分割 $\sigma = (a = x_0 < \dots < x_N = b)$，令
$L_\sigma$ 与 $U_\sigma$ 为阶梯函数，它们在每个
$\intoo{x_{i-1}}{x_i}$ 上分别等于 $m_i = \inf_{[x_{i-1}, x_i]}f$ 与
$M_i = \sup$。达布和即为它们的积分（黎曼积分与勒贝格积分
对阶梯函数一致，二者都给出 $\sum
m_i\Delta x_i$）。取一列分割 $\sigma_n$，
每个都加细前一个，网长 $\to 0$，且达布和
收敛于 $f$ 的下、上达布积分。
加细使 $L_{\sigma_n}$ 不减，$U_{\sigma_n}$ 不增，
在可数集
$D$（全体分点之集）之外逐点成立；记其极限为 $\ell$ 与 $u$
（可测，\cref{prop:b3:lebesgue:stability}）。对 $x \notin
D$，记 $I_n(x)$ 为 $\sigma_n$ 中
含 $x$ 的开区间：$\ell(x) = \sup_n\inf_{I_n(x)}f$ 与 $u(x) =
\inf_n\sup_{I_n(x)}f$。由于网长趋于 $0$，这些
即 $f$ 在 $x$ 处的下、上\emph{包络}
--- $u(x) - \ell(x)$ 即 $f$ 在 $x$ 处的振幅 --- 从而
\emph{$\ell(x) = u(x)$ 当且仅当 $f$ 在 $x$ 处连续}。由单调收敛定理与控制收敛定理（有界，区间有限）：
\[
\int_{\intcc ab}\ell\,\dd\lambda = \lim_n\int L_{\sigma_n}
= \underline{\int}f,
\qquad
\int_{\intcc ab}u\,\dd\lambda = \overline{\int}f .
\]
$f$ 黎曼可积 $\iff$ $\underline\int f =
\overline\int f$ $\iff$ $\int(u - \ell) = 0$ $\iff$ $u = \ell$
几乎处处（$u - \ell \geq 0$；\cref{exo:b3:lebesgue:5}）$\iff$ $f$
几乎处处连续。此时 $\ell \leq f \leq u$，且 $\ell =
u$ 几乎处处：$f$ 几乎处处等于可测函数 $\ell$，因而
为勒贝格可测（由 $\lambda$ 的完备性），且
$\int f\,\dd\lambda = \int\ell\,\dd\lambda = \underline\int f =
\int_a^bf$。
\end{proof}

\begin{example}\label{ex:b3:lebesgue:riemannexamples}
$\mathbf 1_\Q$ 处处不连续：黎曼不可积
--- 但在勒贝格意义下是平凡的：$\int\mathbf 1_\Q\,\dd\lambda =
\lambda(\Q) = 0$。Thomae 函数（\,取 $\frac1q$ 于有理数
$\frac pq$ 处，其余处取 $0$），恰在
无理数处连续：黎曼可积，积分为 $0$。而
\emph{广义}黎曼积分是另一回事：
$\int_0^\infty\frac{\sin x}x\,\dd x$ 作为
$\int_0^A$ 的极限而收敛（周末问题算出其值 $= \frac\pi2$），
但 $\frac{\sin x}x \notin L^1(\intoo0{+\infty})$：绝对
积分如调和级数一样发散
（\cref{exo:b3:lebesgue:6}）。勒贝格理论
以条件收敛换取稳健的极限定理。
\end{example}

\section{含参积分}

以下设 $(X, \mathcal A, \mu)$ 为测度空间，$T$ 为
度量空间（参数），$f \colon T \times X \to \C$，
且 $f(t, \cdot)$ 对每个 $t$ 可积；记 $F(t) = \int_X
f(t, x)\,\dd\mu(x)$。

\begin{theorem}[连续性]\label{thm:b3:lebesgue:paramcont}
假设：$t \mapsto f(t,x)$ 在 $t_0$ 处连续，对几乎处处的
$x$，且存在可积函数 $g$，使得 $\abs{f(t,x)} \leq
g(x)$，对所有 $t$（位于 $t_0$ 的某邻域内）及几乎处处的 $x$ 成立。
则 $F$ 在 $t_0$ 处连续。\index{含参积分}
\end{theorem}

\begin{proof}
对任意数列 $t_n \to t_0$，$f(t_n, \cdot) \to f(t_0,
\cdot)$ 几乎处处收敛，且被 $g$ 控制。由控制收敛定理得 $F(t_n) \to F(t_0)$；
在度量空间中，序列连续性已足够
（\cref{rem:b3:topology:sequences}）。
\end{proof}

\begin{theorem}[积分号下求导]
\label{thm:b3:lebesgue:paramdiff}
设 $T$ 为 $\R$ 中的开区间。假设：对几乎处处的 $x$，
$t \mapsto f(t,x)$ 在 $T$ 上可导，且
\[
\Bigl|\frac{\partial f}{\partial t}(t, x)\Bigr| \leq g(x)
\quad \text{对所有 } t \in T,\ \text{几乎处处 } x,
\]
$g$ 可积。则 $F$ 在 $T$ 上可导，且 $F'(t)
= \int_X \frac{\partial f}{\partial t}(t, x)\,\dd\mu(x)$。
\end{theorem}

\begin{proof}
固定 $t$ 及 $h_n \to 0$，考虑差商
\[
\varphi_n(x) = \frac{f(t + h_n, x) - f(t, x)}{h_n}
\longrightarrow \frac{\partial f}{\partial t}(t,x)
\quad\text{几乎处处},
\]
并由中值不等式得 $\abs{\varphi_n(x)} \leq
\sup_{s}\abs{\partial_tf(s,x)} \leq g(x)$，控制收敛定理适用，且
$\frac{F(t + h_n) - F(t)}{h_n} = \int\varphi_n \to
\int\partial_t f(t, \cdot)$。
\end{proof}

\begin{example}[伽马函数]\label{ex:b3:lebesgue:gamma}
对 $t > 0$，令\index{伽马函数}
\[
\Gamma(t) = \int_0^{+\infty} x^{t-1}\eu^{-x}\,\dd x .
\]
该积分收敛：在 $0$ 附近，$x^{t-1}$ 可积（$t >
0$）；在无穷远处，$x^{t-1}\eu^{-x} \leq C\eu^{-x/2}$。
分部积分（在 $[\varepsilon, A]$ 上，再用
单调收敛定理取极限），得函数方程 $\Gamma(t + 1) =
t\,\Gamma(t)$，由此 $\Gamma(n+1) = n!$，即阶乘的
插值。在每个 $\intcc ab \subseteq \intoo0{+\infty}$ 上，
$\partial_t\bigl(x^{t-1}\eu^{-x}\bigr) = \ln
x\cdot x^{t-1}\eu^{-x}$ 被 $\abs{\ln x}(x^{a-1} +
x^{b-1})\eu^{-x}$ 控制，且可积，于是 $\Gamma$ 属于 $\mathcal C^1$，
由归纳法得 $\mathcal C^\infty$，且 $\Gamma^{(k)}(t) =
\int_0^\infty(\ln x)^kx^{t-1}\eu^{-x}\dd x$。
取值 $\Gamma(\frac12) = \sqrt\pi$ 即是伪装的
高斯积分（\cref{pb:b3:lebesgue:1}）。
\end{example}

\begin{omfigure}
\begin{tikzpicture}
\begin{axis}[omaxis, width=10.6cm, height=5.4cm,
    xmin=0, xmax=25, ymin=-0.25, ymax=1.05,
    xtick={0, 3.1416, 6.2832, 9.4248, 12.566, 15.708, 18.850,
      21.991, 25.133},
    xticklabels={$0$, $\pi$, $2\pi$, $3\pi$, $4\pi$, $5\pi$,
      $6\pi$, $7\pi$, $8\pi$},
    ytick={0, 1}]
  \addplot[omDef!40] coordinates {(0,0) (25,0)};
  \addplot[omThm, thick, domain=0.02:25, samples=600]
    {sin(deg(x))/x};
  \addplot[omDef, dashed, domain=0.02:25, samples=200]
    {1/x};
  \addplot[omDef, dashed, domain=0.02:25, samples=200]
    {-1/x};
\end{axis}
\end{tikzpicture}

{\small 被积函数 $\frac{\sin x}x$：广义积分
$\int_0^\infty$ 收敛，靠的是相邻拱之间的交替抵消，
而各拱的面积 $\abs{\cdot}$ 大致
形如 $\frac2{\pi k}$ --- 即调和级数：$\frac{\sin x}x
\notin L^1$。勒贝格可积即绝对可积。}
\end{omfigure}

\section{练习}

\begin{exercise}[$\star$]\label{exo:b3:lebesgue:1}
(a) 证明单调函数 $\R \to \R$ 是
博雷尔可测的，并证明（处处可导函数的）导函数
是博雷尔可测的。
(b) 证明 $f \colon X \to \R$ 可测，当且仅当 $\{f > q\}
\in \mathcal A$ 对每个\emph{有理数} $q$ 成立。
\end{exercise}

\begin{exercise}[$\star$]\label{exo:b3:lebesgue:2}
计算下列极限，并给出完整论证：
\[
\lim_{n\to\infty}\int_0^{+\infty}
\frac{\cos x}{(1 + x/n)^{n}}\,\dd x,
\qquad
\lim_{n\to\infty}\int_0^1 \frac{n\,x^{n-1}}{1 + x}\,\dd x .
\]
\emph{（对第二个极限：控制之前先作换元 $u = x^n$。）}
\end{exercise}

\begin{exercise}[$\star\star$]\label{exo:b3:lebesgue:3}
(a) 举例说明法图引理中的不等式可以严格成立。
(b) 以三种方式构造逐点有 $f_n \to 0$ 且 $\int f_n = 1$ 的函数列：
即分别在高度上、宽度上与向无穷远处逃逸。
各分别违反控制收敛定理的哪一个假设？
(c) 证明在法图引理中，两侧都不能把 $\liminf$ 换成
$\limsup$。
\end{exercise}

\begin{exercise}[$\star\star$]\label{exo:b3:lebesgue:4}
(a) 证明 $\displaystyle\int_0^{+\infty}\frac{x}{\eu^x -
1}\,\dd x = \sum_{n\geq1}\frac1{n^2} = \frac{\pi^2}6$
\emph{（将 $\frac1{\eu^x - 1}$ 展开为等比级数，并
逐项积分 --- 哪一个定理允许这样做？）}。
(b)（大二学生之梦）证明
$\displaystyle\int_0^1 x^{-x}\,\dd x = \sum_{n\geq1}n^{-n}$。
\emph{（写出 $x^{-x} = \eu^{-x\ln x} = \sum_k\frac{(-x\ln
x)^k}{k!}$，并计算 $\int_0^1(-x\ln x)^k\dd x$，
作换元 $x = \eu^{-u/(k+1)}$，并识别出 $\Gamma$。）}
\end{exercise}

\begin{exercise}[$\star\star$]\label{exo:b3:lebesgue:5}
(a) 证明：若 $f \geq 0$ 可测且 $\int f\,\dd\mu = 0$，
则 $f = 0$ 几乎处处成立。\emph{（考虑 $\{f \geq 1/n\}$ 及
马尔可夫不等式：$\mu(\{f \geq a\}) \leq \frac1a\int f$ ---
请证明之。）}
(b) 证明可积函数 $f$ 几乎处处有限。
(c) 证明：若 $\int_A f\,\dd\mu = 0$ 对\emph{每个}
可测集 $A$，则 $f = 0$ 几乎处处成立。
\end{exercise}

\begin{exercise}[$\star\star$]\label{exo:b3:lebesgue:6}
(a) 运用\cref{thm:b3:lebesgue:riemann}判断下列函数的
黎曼可积性：$\mathbf 1_\Q$；托马埃函数；$\mathbf
1_K$，其中 $K$ 为胖康托集（见\cref{exo:b3:measure:5}）。
(b) 证明 $\int_1^{+\infty}\abs{\frac{\sin x}x}\,\dd x =
+\infty$，而 $\lim_{A\to\infty}\int_1^A\frac{\sin
x}x\,\dd x$ 存在\emph{（用分部积分）}：
广义收敛而不可积。
\end{exercise}

\begin{exercise}[$\star\star$]\label{exo:b3:lebesgue:7}
证明 $F(t) = \int_0^{+\infty}\eu^{-x^2}\cos(tx)\,\dd x$
是 $\mathcal C^1$ 的（定义在 $\R$ 上），且满足 $F'(t) = -\frac t2F(t)$
\emph{（用分部积分）}；由此推出 $F(t) =
F(0)\,\eu^{-t^2/4}$。（由周末问题可知 $F(0) = \frac{\sqrt\pi}2$：
高斯函数本质上就是它自己的傅里叶
变换 --- \cref{ch:b3:fouriertransform} 将
对此进行系统的阐述。）
\end{exercise}

\begin{exercise}[$\star\star\star$]\label{exo:b3:lebesgue:8}
（弗鲁拉尼）设 $0 < a < b$，证明
\[
\int_0^{+\infty}\frac{\eu^{-ax} - \eu^{-bx}}{x}\,\dd x =
\ln\frac ba,
\]
将被积函数写成 $\int_a^b \eu^{-xt}\,\dd t$，并
利用非负函数理论证明积分次序可以交换
（\cref{cor:b3:lebesgue:additivity} 的连续形式 ---
可预先用托内利定理，或将 $[a,b]$ 等分为 $n$ 份，再
取极限）。
\end{exercise}

\begin{exercise}[$\star\star$]\label{exo:b3:lebesgue:9}
设 $f \geq 0$ 为 $(X, \mathcal A, \mu)$ 上的可测函数。证明
$\nu(A) = \int_A f\,\dd\mu$ 定义了一个测度
（\emph{密度}\index{密度} $f$ 关于
$\mu$），并证明 $\int g\,\dd\nu = \int gf\,\dd\mu$
对一切可测的 $g \geq 0$ 成立 \emph{（先对示性函数，再对简单函数，
最后用单调收敛定理——标准方法）}。
\end{exercise}

\begin{exercise}[$\star\star\star$]\label{exo:b3:lebesgue:10}
（一个魏尔斯特拉斯型的反例）。定义 $f(t) =
\int_0^{+\infty}\frac{\sin(tx)}{x(1 + x^2)}\,\dd x$.
(a) 证明 $f$ 在 $\R$ 上有良定义且连续，并
证明它是 $\mathcal C^1$ 的，且 $f'(t) = \int_0^\infty\frac{\cos(tx)}{1 +
x^2}\dd x$ 对每个 $t$ 成立，但
\emph{再次}在积分号下求导是不合法的。
(b) 承认 $f'(t) = \frac\pi2\eu^{-t}$ 对 $t > 0$ 成立（证明
见\cref{ch:b3:residues}）。问
$\int_0^\infty\frac{x\sin(tx)}{1+x^2}\dd x$ 当 $t > 0$ 时等于多少，
以及为什么其公式能印证 (a) 中的失效？
\end{exercise}

\begin{exercise}[$\star\star$]\label{exo:b3:lebesgue:11}
（谢弗引理）。设 $f_n, f \geq 0$ 可积，且
$f_n \to f$ 几乎处处成立，且 $\int f_n \to \int f$。
(a) 证明 $\int\abs{f_n - f} \to 0$。\emph{（
对 $g_n = (f - f_n)^+ \leq f$ 应用控制收敛定理，并
写出 $\int\abs{f_n - f} = 2\int g_n - \int(f - f_n)$。）}
(b) 举例说明假设 $\int f_n \to \int
f$ 不可去掉（滑动或集中的尖峰），并
说明对变号的 $f_n$，结论在缺少
绝对值控制时不成立：$f_n = n\mathbf 1_{\intoc0{1/n}} -
n\mathbf 1_{\intoc{-1/n}0}$ 使 $f_n \to 0$ 几乎处处成立，$\int f_n
= 0 \to 0$，但 $\int\abs{f_n} = 2$。
(c) 应用（密度）：若概率密度 $p_n
\to p$ 几乎处处成立，则自动有 $\int\abs{p_n - p} \to 0$：
密度的逐点收敛即是 $L^1$ 收敛——
无代价的收敛性增强。
\end{exercise}

\begin{exercise}[$\star\star$]\label{exo:b3:lebesgue:12}
经典极限，借助单调收敛定理与控制收敛定理给出完整论证：
\[
\text{(a)}\ \lim_{n\to\infty}\int_0^n\Bigl(1 -
\frac xn\Bigr)^n\eu^{x/2}\,\dd x, \qquad
\text{(b)}\ \lim_{n\to\infty}\int_0^1\frac{n\,x^{n-1}}{1 +
x}\,\dd x,
\]
\[
\text{(c)}\ \lim_{n\to\infty}\int_0^\infty
\frac{\dd x}{(1 + x/n)^n\,x^{1/n}} .
\]
\emph{（对 (a)：$(1 - x/n)^n \nearrow \eu^{-x}$ 对固定的
$x$ 成立，单调性借助 $\log$ 证明；对 (b)，
分部积分，或作代换 $x = u^{1/n}$，并识别
边界处的集中现象；对 (c)，找一个可积
对所有 $n \geq 2$ 都成立的控制函数，办法是在 $x =
1$ 处分段。）}
\end{exercise}

\section{周末问题：两个著名积分}

\begin{problem}[{周末问题 --- 高斯积分与
狄利克雷积分，仅用参数方法}]
\label{pb:b3:lebesgue:1}
两个积分在应用分析中占据核心地位：
\[
G = \int_{-\infty}^{+\infty}\eu^{-x^2}\dd x = \sqrt\pi,
\qquad
D = \int_0^{+\infty}\frac{\sin x}{x}\,\dd x = \frac\pi2
\]
（第二个积分按反常积分理解，
见\cref{ex:b3:lebesgue:riemannexamples}）。仅用本章
的工具即可证明两者。

\medskip
\noindent\textbf{第一部分 --- 高斯积分。}对 $t \geq 0$，记
\[
A(t) = \Bigl(\int_0^t\eu^{-x^2}\dd x\Bigr)^{2},
\qquad
B(t) = \int_0^1\frac{\eu^{-t^2(1 + x^2)}}{1 + x^2}\,\dd x .
\]
\begin{enumerate}
  \item 证明 $A$ 与 $B$ 为 $\mathcal C^1$ 函数，定义域为
        $\intoo0{+\infty}$，并求 $A'$ 与 $B'$；证明
        $A'(t) + B'(t) = 0$。\emph{（求 $B'$ 时作代换 $u =
        tx$。）}
  \item 计算 $A(0) + B(0)$ 与 $\lim_{t\to+\infty}(A + B)(t)$，
        ——对 $B$ 中的积分，证明极限可移入积分号。
  \item 由此得 $\int_0^\infty \eu^{-x^2}\dd x =
        \frac{\sqrt\pi}2$，从而 $G = \sqrt\pi$，并推出
        $\Gamma(\tfrac12) = \sqrt\pi$ \emph{（作代换 $x =
        u^2$，并代入 $\Gamma(\frac12)$）}。
\end{enumerate}

\medskip
\noindent\textbf{第二部分 --- 狄利克雷积分。}对 $t
\geq 0$，记
\[
F(t) = \int_0^{+\infty}\eu^{-tx}\,\frac{\sin x}{x}\,\dd x .
\]
\begin{enumerate}[resume]
  \item 证明定义 $F(t)$ 的积分对每个
        $t > 0$ 在勒贝格积分意义下收敛，对 $t = 0$
        作为反常积分收敛；证明 $D =
        \lim_{A\to\infty}\int_0^A\frac{\sin x}x\dd x$ 存在
        \emph{（在 $[\pi, A]$ 上分部积分）}。
  \item 证明 $F$ 为 $\mathcal C^1$ 函数，定义域为
        $\intoo0{+\infty}$，且
        \[
        F'(t) = -\int_0^{+\infty}\eu^{-tx}\sin x\,\dd x
        = -\frac{1}{1 + t^2}
        \]
        \emph{（在 $[t_0, \infty)$ 上控制，对每个 $t_0 >
        0$；最后一个积分用两次分部积分，或
        用复指数）}。
  \item 证明 $F(t) \to 0$（$t \to +\infty$），并由此推出 $F(t)
        = \frac\pi2 - \arctan t$ 在 $\intoo0{+\infty}$ 上成立。
  \item 难点在于 $D = \lim_{t\to0^+}F(t)$。证明时
        以尾部一致控制完成：对 $0 \leq t \leq 1$
        及 $A \geq \pi$，分部积分以证明
        \[
        \Bigl|\int_A^{+\infty}\eu^{-tx}\frac{\sin
        x}x\,\dd x\Bigr| \leq \frac{C}{A}
        \]
        其中 $C$ 与 $t$ 无关。\emph{（对
        $\frac{\eu^{-tx}}x$ 求导，并将 $\abs{\cos}$ 以 $1$ 为界；
        注意 $t\eu^{-tx} \leq 1/x\cdot(tx\eu^{-tx})$，并利用
        $\sup_{u\geq0}u\eu^{-u} < 1$）}；随后把 $F(t) - D$
        分拆为 $[0, A]$（其上当 $t \to 0$ 时可用控制收敛定理）与
        $[A, \infty)$ 两部分。
  \item 综上，$D = \frac\pi2$。
\end{enumerate}

\medskip
\noindent\textbf{第三部分 --- 红利。}
\begin{enumerate}[resume]
  \item 计算 $\int_0^{+\infty}\frac{\sin^2x}{x^2}\,\dd x$
        \emph{（分部积分，化为 $D$，并用到 $\sin
        2x = 2\sin x\cos x$）}。
  \item 计算 $\int_0^{+\infty}\frac{1 -
        \cos x}{x^2}\,\dd x$，并检验
        两个结果是否一致。
  \item 设 $a > 0$，计算
        $\int_0^{+\infty}\frac{\sin(ax)}x\dd x$ 与
        $\int_{-\infty}^{+\infty}\eu^{-ax^2}\dd x$，并记下
        尺度变换规律（它们将在
        \cref{ch:b3:fouriertransform}中发挥主要作用）。
  \item 精确说明为何 $D$ 不能
        在 $t = 0$ 处直接用控制收敛定理处理（
        $[0,1]\times[0,\infty)$ 上没有可积的控制函数），以及
        为何第 7 题中的尾部分拆是恰当的替代办法 --- 这一模式
        （“尾部一致可积性”）在整个
        分析学中反复出现。
\end{enumerate}

\medskip
\noindent\textbf{第四部分 --- 依
玻尔--莫勒鲁普的伽马函数。} 函数 $\Gamma$
        （\cref{ex:b3:lebesgue:gamma}）满足 $\Gamma(1) = 1$ 及
        $\Gamma(x+1) = x\Gamma(x)$ --- 但满足它们的函数还有
        无穷多个（只需乘以任一周期为 $1$ 的振荡函数）。一个凸性
        条件即可唯一确定 $\Gamma$，其更深的
        恒等式便随之机械地导出。称区间上的正函数 $f$
        为\emph{对数凸}，若 $\log f$ 为凸函数。
\begin{enumerate}[resume]
  \item 证明对数凸蕴含凸，对数凸函数之积
        及其与仿射映射的复合仍为对数凸，并借助
        双函数形式的
        赫尔德不等式 $\int\abs{uv} \leq
        \bigl(\int\abs u^p\bigr)^{1/p}\bigl(\int\abs
        v^q\bigr)^{1/q}$（由杨氏不等式直接证明） ---
        证明 $\Gamma$ 在
        $\intoo0\infty$ 上对数凸。
  \item （斜率引理）设 $g$ 为 $\intoo0\infty$ 上的凸函数，
        且 $g(n+1) - g(n) = \log n$，对每个整数 $n
        \geq 1$。对 $x \in \intoc01$ 及 $n \geq 2$，比较
        $g$ 在 $[n-1, n]$、$[n, n+x]$ 与
        $[n, n+1]$ 上的斜率，并推出
        \[
        x\log(n-1) \;\leq\; g(n + x) - g(n) \;\leq\; x\log n .
        \]
  \item （玻尔--莫勒鲁普）设 $f > 0$ 满足 $f(1) = 1$，
        $f(x+1) = xf(x)$，且 $\log f$ 为凸函数。将递推式
        展开为 $f(n + x) = x(x+1)\cdots(x + n -
        1)\,f(x)$ 及 $f(n) = (n-1)!$，由第 14 题
        可知对 $x \in \intoc01$ 有
        \[
        f(x) = \lim_{n\to\infty}
        \frac{n!\,n^x}{x(x+1)\cdots(x+n)} :
        \]
        $f$ 唯一，故 $f = \Gamma$，且\emph{高斯
        极限公式}成立（借助递推式可推广至所有 $x > 0$
        的情形）。
  \item 定义\emph{贝塔函数} $B(x, y) =
        \int_0^1t^{x-1}(1-t)^{y-1}\,\dd t$（$x, y > 0$）。证明
        收敛性，以及递推式 $B(x+1, y) =
        \frac{x}{x+y}\,B(x, y)$（用分部积分）及
        $B(1, y) = \frac1y$。
  \item 证明 $x \mapsto B(x, y)$ 对数凸（
        再次用到赫尔德不等式），并将玻尔--莫勒鲁普定理用于
        \[
        f(x) = \frac{B(x, y)\,\Gamma(x + y)}{\Gamma(y)}
        \]
        从而得出\emph{欧拉公式}：$B(x, y) =
        \dfrac{\Gamma(x)\Gamma(y)}{\Gamma(x+y)}$ --- 全程不涉及
        任何二重积分。
  \item 直接计算 $B(\frac12, \frac12)$（作代换 $t
        = \sin^2\theta$），并由此推得 $\Gamma(\frac12) =
        \sqrt\pi$：即第一部分中的高斯积分，此处仅由
        凸性即可恢复。将两种证明各用一句话
        加以比较。
  \item （勒让德倍元公式）证明
        \[
        g(x) = \frac{2^{x-1}}{\sqrt\pi}\,
        \Gamma\Bigl(\frac x2\Bigr)
        \Gamma\Bigl(\frac{x+1}2\Bigr)
        \]
        满足玻尔--莫勒鲁普的三个假设，并
        由此得 $g = \Gamma$，即 $\Gamma(2z) =
        \frac{2^{2z-1}}{\sqrt\pi}\,\Gamma(z)\,\Gamma(z +
        \tfrac12)$ 对所有 $z > 0$ 成立。
  \item 由此推出闭形式 $\Gamma\bigl(n + \tfrac12\bigr)
        = \dfrac{(2n)!}{4^n\,n!}\sqrt\pi$，并
        借助斜率引理，对 $\log\Gamma$ 在
        大整数附近的情形证明下述渐近式
        \[
        \frac{\Gamma(n + \frac12)}{\Gamma(n)\,\sqrt n}
        \longrightarrow 1 .
        \]
  \item 将最后两题合并，得到中心
        二项式渐近式
        \[
        \binom{2n}{n} \sim \frac{4^n}{\sqrt{\pi n}},
        \]
        并对 $n = 10$ 作数值验证（$\binom{20}{10} =
        184756$，对照 $4^{10}/\sqrt{10\pi} \approx
        187079$：比值 $\approx 0.988$）。
  \item （综合）常数 $\sqrt\pi$ 已先后以
        高斯积分（第一部分）、$B(\frac12,
        \frac12)$（第 18 题）出现，并
        在倍元公式中出现（第 19 题）；中心二项式估计
        预示了斯特林公式（\cref{ch:b3:product}中的
        周末问题）与棣莫弗--拉普拉斯定理。梳理
        这些联系：哪些陈述彼此等价，
        以及积分号下求导与凸性两种技巧 ---
        各自贡献了什么，是对方所不能提供的
        内容？
\end{enumerate}

\medskip
\noindent\textbf{Part V --- 另外三项收获。}
\begin{enumerate}[resume]
  \item （沃利斯公式，借助贝塔函数）设 $p > -1$，作代换 $t =
        \sin^2\theta$，证明
        \[
        W_p = \int_0^{\pi/2}\sin^p\theta\,\dd\theta
        = \frac12\,B\Bigl(\frac{p+1}2, \frac12\Bigr),
        \]
        由贝塔递推关系（第 16 题）可得
        $W_{n+2} = \frac{n+1}{n+2}\,W_n$，对整数 $n \geq
        0$。求 $W_{2n}$ 与 $W_{2n+1}$ 的闭形式，
        用夹逼法证明 $W_{2n+1}/W_{2n} \to 1$，并借此得到
        \emph{沃利斯公式}
        \[
        \frac\pi2 = \lim_{n\to\infty}
        \prod_{k=1}^{n}\frac{4k^2}{4k^2 - 1} .
        \]
  \item （高斯函数遇见频率）对 $b \in \R$，令
        \[
        \Phi(b) = \int_{-\infty}^{+\infty}
        \eu^{-x^2}\cos(2bx)\,\dd x .
        \]
        证明 $\Phi$ 属于 $\mathcal C^1$（在 $\R$ 上），
        由分部积分得到微分方程
        $\Phi'(b) = -2b\,\Phi(b)$，并得出
        \[
        \Phi(b) = \sqrt\pi\,\eu^{-b^2} :
        \]
        高斯函数在此变换下自我再现 ---
        这是唯一的恒等式，
        \cref{ch:b3:fouriertransform}正是在此基础上展开的。
  \item （弗鲁拉尼积分）设 $0 < a < b$，证明
        \[
        \int_0^{+\infty}
        \frac{\eu^{-ax} - \eu^{-bx}}{x}\,\dd x
        = \log\frac ba ,
        \]
        通过对参数 $a$ 求导（说明
        对每个 $\intco{a_0}{+\infty}$、$a_0 > 0$ 的控制条件成立），
        并令 $a \to b$ 以确定常数）。
        被积函数究竟在何处需要处理其可去奇点，
        即 $x = 0$ 处的奇点？
\end{enumerate}
\end{problem}
